\section{Lei de Gauss (eletricidade)}

A lei de Gauss relaciona o campo elétrico com a distribuição de carga. Dada uma superfície orientada $S$, e um campo vetorial $\mathbf F$, definimos o fluxo de $\mathbf{F}$ através de $S$ como sendo $\Phi_F(S) = \oiint_S \mathbf{F} \cdot d\mathbf{S}$.

Para o campo elétrico $\mathbf{E}$, a lei de Gauss afirma que o fluxo de $\mathbf{E}$ através de uma superfície fechada orientada (gaussiana) $S$ é proporcional à carga total localizada em seu interior, $Q$:
\begin{equation}
\oiint_{S} \mathbf{E} \cdot d\mathbf{S} = Q / \epsilon_0 ,
\label{eq: gauss_electric}
\end{equation} 

\noindent onde $\epsilon_0$ é a permissividade elétrica do vácuo. Considere uma carga pontual $q$ e uma esfera de raio $r$ centrada na carga. O campo devido à carga é perpendicular à superfície da esfera em todos os seus pontos, e seu módulo depende apenas de $r$, de modo que $\mathbf{E} \cdot d\mathbf{S}$ é uma constante. Aplicando a eq. \eqref{eq: gauss_electric} e isolando $\mathbf{E}$, obtemos a expressão para o campo elétrico devido a uma carga pontual, que concorda com a expressão obtida através da lei de Coulomb:
\begin{equation}
\mathbf{E} = \frac{q}{4 \pi r^2 \epsilon_0}.
\label{eq: coulomb} 
\end{equation}

Fazendo uso do teorema da divergência \eqref{th: div}, podemos reescrever a eq. \eqref{eq: gauss_electric} como
\begin{equation}
\frac{Q}{\epsilon_0}
= \oiint_S \mathbf{E} \cdot d\mathbf{S}
= \iiint_V \nabla \cdot \mathbf{E} \,dV.
\label{eq: div_theorem_gauss}
\end{equation}

\noindent Se $\rho$ é a densidade volumétrica de carga, temos que $Q = \iiint_V \rho \,dV$. Substituindo em \eqref{eq: div_theorem_gauss}, obtemos a forma diferencial da lei de Gauss, também conhecida como equação de Poisson:
\begin{equation}
\nabla \cdot \mathbf{E} = \rho / \epsilon_0.
\label{eq: poisson}
\end{equation}

